\documentclass[11pt,reqno]{amsart}
\usepackage[T1]{fontenc}
\usepackage{lmodern}
\usepackage[a4paper,margin=27mm]{geometry}
\usepackage{microtype,amsmath,amssymb,mathtools}
\usepackage[hidelinks,pdfusetitle]{hyperref}
\usepackage{enumitem,booktabs,needspace}
\numberwithin{equation}{section}
\newtheorem{theorem}{Theorem}[section]
\newtheorem{proposition}[theorem]{Proposition}
\newtheorem{lemma}[theorem]{Lemma}
\newtheorem{corollary}[theorem]{Corollary}
\theoremstyle{definition}
\newtheorem{definition}[theorem]{Definition}
\newtheorem{remark}[theorem]{Remark}
\newtheorem{example}[theorem]{Example}
\newcommand{\R}{\mathbb R}
\newcommand{\Z}{\mathbb Z}
\newcommand{\N}{\mathbb N}
\newcommand{\E}{\mathbb E}
\newcommand{\Prob}{\mathbb P}
\newcommand{\1}{\mathbf 1}
\newcommand{\dd}{\,d}
\newcommand{\abs}[1]{\left|#1\right|}
\newcommand{\norm}[1]{\left\lVert#1\right\rVert}
\DeclareMathOperator{\Var}{Var}
\DeclareMathOperator{\Ent}{Ent}
\DeclareMathOperator{\PV}{PV}
\DeclareMathOperator{\sgn}{sgn}
\DeclareMathOperator{\supp}{supp}
\DeclareMathOperator{\diver}{div}
\DeclareMathOperator{\Conf}{Conf}
\allowdisplaybreaks[2]
\setlength{\emergencystretch}{1.5em}
\setlist[enumerate]{leftmargin=*,itemsep=3pt,topsep=5pt}
\raggedbottom
\title[Discrepancy-based DLR uniqueness]{Discrepancy-based uniqueness for\\
one-dimensional repulsive Gibbs states}
\date{Research note --- 14 September 2026}
\begin{document}
\begin{abstract}
We extend a logarithmic DLR comparison argument to inverse-power
interactions and to smooth finite-range perturbations of the logarithm.
For each $s>0$ and $\beta>0$, there is at most one stationary
intensity-one canonical DLR state for $r^{-s}$ whose interval-count
variance is bounded linearly and is $o(t)$. An entropy comparison
between nearby temperatures compensates for the change of coupling
under affine normalization; a weighted half-line estimate makes the
exterior-force cost tight. For $-a\log r+W(r)$, with $W$ smooth and of
finite range and the total interaction having positive second
derivative, uniqueness follows under the logarithmic discrepancy
assumptions and a finite local third count moment. We include the
probabilistic and analytic comparison proofs, distinguish these
statements from finite-energy uniqueness, and give the tangential
Hessian obstruction to the same proof in higher dimensions.
\end{abstract}
\maketitle

\noindent\textit{Scope.}
This note develops the comparison framework of the supplied working
manuscript~\cite{Draft}. All model-specific estimates needed below are
proved here. Standard convexity inequalities, Campbell's formula, the
point-shift invariance of Palm measure, and the ergodic theorem are used
in their usual forms. We prove uniqueness within the stated classes;
we do not prove their nonemptiness or derive their count bounds from
finite energy. No claim of priority for the uniqueness statements is made.

\section{Statements and the canonical specification}
\label{sec:statements}

Let $\Conf(\R)$ be the space of locally finite simple counting measures
on $\R$. Write $\theta_a\gamma=\gamma-a$ and $N_I(\gamma)=\gamma(I)$.
A stationary law $P$ has intensity one if
$\E_PN_I=|I|$ for every bounded interval $I$. Put
\[
v_P(t)=\Var_PN_{(0,t]},\qquad t>0.
\]
Throughout, stationarity means invariance under every real translation;
rotation or reflection invariance is not imposed. Constants may depend
on the candidate process.

\begin{theorem}[Inverse-power interactions]\label{thm:riesz}
Fix $s>0$ and $\beta>0$. There is at most one stationary simple
intensity-one canonical DLR state for
\[
V_s(r)=r^{-s},\qquad r>0,
\]
satisfying
\begin{align}
v_P(t)&\le C_Pt\quad(t>0),\label{hyp:linear}\\
v_P(t)&=o(t)\quad(t\longrightarrow\infty).\label{hyp:sublinear}
\end{align}
The DLR specification is the relative-energy specification defined below.
\end{theorem}

\begin{theorem}[Smooth perturbations of the logarithm]\label{thm:perturb}
Let $a_0>0$ and let $W\in C^2([0,\infty))$ vanish on $[R_0,\infty)$
for some $R_0<\infty$. Suppose
\begin{equation}\label{hyp:convex}
V(r)=-a_0\log r+W(r),\qquad V''(r)>0\quad(r>0).
\end{equation}
For each $\beta>0$, there is at most one stationary simple
intensity-one canonical DLR state for $V$ satisfying
\eqref{hyp:linear}--\eqref{hyp:sublinear} and
\begin{align}
\int_1^\infty v_P(t)t^{-2}\dd t&<\infty,\label{hyp:integrated}\\
\E_PN_{[0,1]}^3&<\infty.\label{hyp:third}
\end{align}
\end{theorem}

\begin{example}\label{ex:compact}
For $a_0,R_0>0$ and $\lambda\ge0$,
\[
V(r)=-a_0\log r+\lambda(1-r/R_0)_+^4
\]
satisfies \eqref{hyp:convex}. Indeed, the added second derivative is
$12\lambda R_0^{-2}(1-r/R_0)^2$ for $r<R_0$ and zero for $r>R_0$,
with continuous extension at $R_0$. No smallness condition on $\lambda$
is required. The count hypotheses remain assumptions on the Gibbs state.
\end{example}

\subsection{Relative exterior energy}
Let $I=(A,B)$ be bounded, $\xi$ an exterior configuration supported on
$I^c$, and $x_*\in I$. For $x\in I$, define
\begin{equation}\label{def:exterior}
U_{V,I}^{\xi}(x)=
\lim_{R\to\infty}\sum_{\substack{p\in\xi\\ |p|\le R}}
\bigl[V(|x-p|)-V(|x_*-p|)\bigr],
\end{equation}
whenever the limit exists. For $V_s$ it is an absolutely convergent
series on density-one exteriors. In fact, uniformly on compact sets of
$x,x_*$,
\begin{equation}\label{eq:relative-riesz-tail}
\abs{V_s(|x-p|)-V_s(|x_*-p|)}
\le C_{I,s}|p|^{-s-1}\quad\hbox{for large }|p|.
\end{equation}
For a logarithmic tail the cutoff is symmetric. The discrepancy
assumption \eqref{hyp:integrated} gives the reciprocal principal value
needed for convergence; this is proved in
Section~\ref{sec:forces}. The finite-range contribution is a finite sum.
Changing $x_*$ adds an $x$-independent constant.

On the ordered chamber
$I_<^n=\{A<x_1<\cdots<x_n<B\}$, let
\begin{equation}\label{def:hamiltonian}
H_{V,I}^{\xi,n}(x_1,\ldots,x_n)
=\sum_{1\le i<j\le n}V(x_j-x_i)
+\sum_{i=1}^nU_{V,I}^{\xi}(x_i).
\end{equation}
The associated canonical probability is
\begin{equation}\label{def:kernel}
\pi_{V,I,n}^{\xi}(\dd x)=
\frac{\1_{I_<^n}(x)}{Z_{V,I,n}^{\xi}}
\exp\{-\beta H_{V,I}^{\xi,n}(x)\}\dd x_1\cdots\dd x_n.
\end{equation}
For $n=0$ it is the point mass at the empty configuration. Ordering
absorbs the usual $n!$ factor into the normalization.

\begin{definition}[Canonical DLR equation]\label{def:dlr}
A point-process law $P$ satisfies the canonical DLR equation for $V$
at inverse temperature $\beta$ if, for every bounded interval $I$,
conditionally on $(\gamma_{I^c},N_I)$ its interior law is
$\pi_{V,I,N_I}^{\gamma_{I^c}}$ almost surely. Equivalently, for every
bounded measurable $f$,
\begin{equation}\label{eq:dlr}
\E_P\!\left[f(\gamma)\mid\gamma_{I^c},N_I\right]
=\int f\!\left(\gamma_{I^c}+\sum_{i=1}^{N_I}\delta_{x_i}\right)
\pi_{V,I,N_I}^{\gamma_{I^c}}(\dd x).
\end{equation}
The relative energies and normalizations in this definition are
required to be well defined on almost every exterior.
\end{definition}

For $0<s<1$, equations \eqref{def:exterior}--\eqref{eq:dlr} in particular
give the canonical Riesz DLR equation without an infinite unrenormalized
exterior energy. The anchor subtraction is essential at the level of
that energy, even though its force is absolutely summable.
For the logarithmic interaction this convention is the canonical
relative-energy convention considered in~\cite{DHLM}.

\subsection{The lattice reference}
For $K\ge1$, set $L=K+1$ and
\[
\Omega_K=\{0<u_1<\cdots<u_K<L\},\qquad
\xi_K^{\mathrm{lat}}=\{j\in\Z:j\le0\text{ or }j\ge L\}.
\]
Let
\begin{equation}\label{def:lattice-H}
H^{\mathrm{lat}}_{V,K}(u)=
\sum_{i<j}V(u_j-u_i)+
\sum_i\PV\sum_{p\in\xi_K^{\mathrm{lat}}}
\bigl[V(|u_i-p|)-V(|x_*-p|)\bigr].
\end{equation}
For Riesz interactions the $\PV$ sign may be dropped.
For the logarithm, pair $p=-m$ with $p=L+m$: the relative logarithms
have a paired tail $O(m^{-2})$, locally uniformly in $u$.
The cutoff convention in \eqref{def:lattice-H} agrees with this pairing,
since the finitely shifted cutoff shells contribute $O(R^{-1})$.
Derivatives of order two converge locally uniformly in both models.

Every pair Hessian is nonnegative in $\Omega_K$, and the endpoint
terms are convex as well. Thus $H^{\mathrm{lat}}_{V,K}$ is convex and
$C^2$ there. A supporting affine function bounds it below on the
bounded chamber. Consequently, for every $T,\tau>0$,
\begin{align}
\Phi_{K,T}(u)&=\beta H^{\mathrm{lat}}_{V,K}(u)
+\frac1{2T}\sum_{i=1}^K(u_i-i)^2,\label{def:Phi}\\
\omega_{K,T}^{(\tau)}(\dd u)
&=Z_{K,T}(\tau)^{-1}e^{-\tau\Phi_{K,T}(u)}\dd u
\label{def:omega-tau}
\end{align}
have strictly positive finite normalizing constants. We write
$\omega_{K,T}=\omega_{K,T}^{(1)}$.


\section{A static inequality for convex one-dimensional interactions}
\label{sta:sec-static}\label{sec:static}

Fix $\beta>0$. Let $V\in C^2((0,\infty))$ satisfy
\begin{equation}\label{sta:curvature}
V''(t)>0\qquad(t>0).
\end{equation}
This is a positive-curvature assumption; strict convexity alone would
not imply it. Write $L=K+1$ and
\[
\Omega_K=\{0<u_1<\cdots<u_K<L\},\qquad
u_0=0,\quad u_{K+1}=L,\quad g_i=u_{i+1}-u_i.
\]
The lattice exterior is $\{j\in\mathbb Z:j\le0\text{ or }j\ge L\}$.
Assume its relative Hamiltonian $H^{\mathrm{lat}}_{V,K}$ is finite and
$C^2$ on $\Omega_K$, and that its Hessian is the sum of the
nonnegative Hessians of the interior pairs and exterior interactions.
In particular, for $v_0=v_{K+1}=0$,
\begin{equation}\label{sta:hessian}
v^\top D^2(\beta H^{\mathrm{lat}}_{V,K})(u)v
\ge \beta\sum_{i=0}^K V''(g_i)(v_{i+1}-v_i)^2.
\end{equation}
The nearest-neighbour interior pairs and the two exterior charges at
$0,L$ give the terms on the right. The existence and differentiability
of the relative Hamiltonian must be verified for the interaction under
consideration; they do not follow from \eqref{sta:curvature} alone.

For $T>0$, set
\begin{equation}\label{sta:reference}
\Phi_{K,T}(u)=\beta H^{\mathrm{lat}}_{V,K}(u)
 +\frac1{2T}\sum_{i=1}^K(u_i-i)^2,
\qquad
\omega_{K,T}(\dd u)=Z_{K,T}^{-1}e^{-\Phi_{K,T}(u)}\dd u.
\end{equation}
A finite convex function on this bounded chamber has a global affine
lower bound, obtained from any supporting tangent plane. Its exponential
therefore has a finite integral, and the integral is strictly positive.
This justifies $0<Z_{K,T}<\infty$ under the preceding assumptions.

Fix $r\ge1$ and $G\in C_c^2((0,\infty)^r)$, and write
\begin{equation}\label{sta:observable}
F_K(u)=\frac1{K-r}\sum_{j=1}^{K-r}
 G(g_j,\ldots,g_{j+r-1}),\qquad K\ge2r+1.
\end{equation}
For $\mu\ll\omega$, use the conventions
\[
S(\mu\mid\omega)=\int\log\frac{\dd\mu}{\dd\omega}\,\dd\mu,
\qquad
I(\mu\mid\omega)=4\int
 \left|\nabla\sqrt{\frac{\dd\mu}{\dd\omega}}\right|^2\dd\omega.
\]
The last expression defines Fisher information in the weak Sobolev
sense and is assigned the value $+\infty$ outside its domain. For a
positive smooth density it equals the usual integral of its squared
logarithmic gradient against $\mu$.

\begin{proposition}[Uniform gap comparison]\label{sta:prop-gap}
There is a constant $C$, depending on $V,\beta,r,G$ but independent of
$K\ge2r+1$ and $T>0$, such that
\begin{align}
|\mu(F_K)-\omega_{K,T}(F_K)|^2
 &\le \frac{C}{K}S(\mu\mid\omega_{K,T}),
 \label{sta:entropy-gap}\\
|\mu(F_K)-\omega_{K,T}(F_K)|^2
 &\le \frac{CT}{K}I(\mu\mid\omega_{K,T}).
 \label{sta:fisher-gap}
\end{align}
The same constants can be chosen uniformly when $\beta$ varies in a
fixed compact subinterval of $(0,\infty)$.
\end{proposition}

\begin{proof}
Write $\Phi=\Phi_{K,T}$ and $n=K-r$. Regard $F_K$ as a function of
the gaps and define $q_i=\partial F_K/\partial g_i$, including
$q_0=q_K=0$. Its Euclidean gradient is explicitly
\begin{equation}\label{sta:gradient}
(\nabla F_K)_j=q_{j-1}-q_j\quad(1\le j\le K),
\qquad
\nabla F_K\cdot v=\sum_{i=0}^Kq_i(v_{i+1}-v_i).
\end{equation}
Each gap occurs in at most $r$ summands of \eqref{sta:observable}.
Consequently $|q_i|\le r\|DG\|_\infty/n$. There is a compact interval
$[b,B]\subset(0,\infty)$ containing every coordinate projection of
the supports of $DG$ and $D^2G$. Whenever $q_i\ne0$, one has
$g_i\in[b,B]$. Put $m=\min_{[b,B]}V''>0$.
Using $n\ge K/2$ and \eqref{sta:hessian}, weighted Cauchy--Schwarz gives
\begin{align}
|\nabla F_K\cdot v|^2
&\le\left(\sum_{i:q_i\ne0}
 \frac{q_i^2}{\beta V''(g_i)}\right)
 \left(\beta\sum_{i=0}^K V''(g_i)(v_{i+1}-v_i)^2\right)
 \notag\\
&\le\frac{A}{K}\,v^\top D^2\Phi\,v.
\label{sta:metric-form}
\end{align}
Since $D^2\Phi\ge T^{-1}\mathrm{Id}$, duality yields
\begin{equation}\label{sta:metric-gradient}
\nabla F_K^\top(D^2\Phi)^{-1}\nabla F_K\le A/K.
\end{equation}

There is also a constant $B_0>0$ such that
\begin{equation}\label{sta:tilt-control}
|v^\top D^2F_K\,v|
\le\frac{B_0}{K}\,
 v^\top D^2(\beta H^{\mathrm{lat}}_{V,K})\,v.
\end{equation}
Indeed, for the $j$th window let
$w_j=(v_{j+1}-v_j,\ldots,v_{j+r}-v_{j+r-1})$. Its contribution to the
left is bounded by $n^{-1}\|D^2G\|_\infty|w_j|^2$, and vanishes
unless all its relevant gaps belong to $[b,B]$. Summing, using the
overlap bound $r$ and then the lower bound $m$ in
\eqref{sta:hessian}, proves \eqref{sta:tilt-control}.

For a real parameter $z$, let
\[
\omega_z(\dd u)=
 \frac{e^{zF_K(u)}}{\omega_{K,T}(e^{zF_K})}\omega_{K,T}(\dd u).
\]
Choose $c>0$ with $cB_0\le1/2$. For $|z|\le cK$,
\begin{equation}\label{sta:tilt-convexity}
D^2(\Phi-zF_K)
\ge\tfrac12D^2(\beta H^{\mathrm{lat}}_{V,K})
   +T^{-1}\mathrm{Id}
\ge\tfrac12D^2\Phi.
\end{equation}
The Brascamp--Lieb variance inequality on a convex domain thus implies
\begin{equation}\label{sta:tilted-variance}
\Var_{\omega_z}(F_K)
\le\int\nabla F_K^\top[D^2(\Phi-zF_K)]^{-1}\nabla F_K\,\dd\omega_z
\le 2A/K.
\end{equation}
For completeness, the convex-domain version follows from the Neumann
Bochner identity. For $\nu\propto e^{-\Psi}$ on a smooth convex bounded
domain, let $\mathcal L=\Delta-\nabla\Psi\cdot\nabla$ and solve
$-\mathcal Lh=f-\nu(f)$ with Neumann boundary condition. Integration
by parts and the nonnegative second fundamental form of the boundary
give
\[
\int(\mathcal Lh)^2\dd\nu
\ge\int\nabla h^\top D^2\Psi\,\nabla h\,\dd\nu,
\quad
\Var_\nu(f)=\int\nabla f\cdot\nabla h\,\dd\nu.
\]
Cauchy--Schwarz in the $D^2\Psi$ metric gives
\[
\Var_\nu(f)^2\le
 \left(\int\nabla f^\top(D^2\Psi)^{-1}\nabla f\,\dd\nu\right)
 \Var_\nu(f),
\]
which is the asserted inequality. Smooth convex approximation treats
polytopes. To handle the singularities of $\Phi$ at the original
chamber boundary, first apply this inequality on the inner convex
polytopes
\[
\Omega_{K,\varepsilon}
 =\{u\in\Omega_K:g_i(u)>\varepsilon\text{ for }0\le i\le K\},
\qquad 0<\varepsilon<1.
\]
All derivatives of the potential are regular on their closures. The
constants in \eqref{sta:metric-gradient} and
\eqref{sta:tilt-convexity} are unchanged. The normalized restrictions
converge to $\omega_z$, and $F_K$ is bounded; hence their variance
bounds imply \eqref{sta:tilted-variance} on $\Omega_K$.

Let
$\Lambda(z)=\log\omega_{K,T}
(e^{z(F_K-\omega_{K,T}(F_K))})$.
Boundedness of $F_K$ permits differentiation under the integral.
Since $\Lambda(0)=\Lambda'(0)=0$ and
$\Lambda''(z)=\Var_{\omega_z}(F_K)$, integrating
\eqref{sta:tilted-variance} twice, for either sign of $z$, gives
\begin{equation}\label{sta:mgf}
\Lambda(z)\le A z^2/K\qquad(|z|\le cK).
\end{equation}
The Donsker--Varadhan entropy variational inequality, applied to both $zF_K$ and
$-zF_K$, now shows, for $0<z\le cK$,
\begin{equation}\label{sta:entropy-optimization}
|\mu(F_K)-\omega_{K,T}(F_K)|
\le S(\mu\mid\omega_{K,T})/z+Az/K.
\end{equation}
If $0<S\le Ac^2K$, choose $z=\sqrt{SK/A}$ to obtain a squared bound
$4AS/K$; the case $S=0$ follows by taking $z\downarrow0$.
If $S>Ac^2K$, boundedness gives instead
\[
|\mu(F_K)-\omega_{K,T}(F_K)|^2
\le4\|G\|_\infty^2
\le \frac{4\|G\|_\infty^2}{Ac^2}\frac SK.
\]
The infinite-entropy case is automatic. This proves
\eqref{sta:entropy-gap} for every entropy value.

Finally $D^2\Phi\ge T^{-1}\mathrm{Id}$ gives the convex-domain
logarithmic Sobolev inequality
\begin{equation}\label{sta:lsi}
\Ent_{\omega_{K,T}}(f^2)
\le2T\int|\nabla f|^2\dd\omega_{K,T},
\qquad
S(\mu\mid\omega_{K,T})\le\frac T2 I(\mu\mid\omega_{K,T}).
\end{equation}
Here $\Ent_\omega(h)=\omega(h\log h)-\omega(h)\log\omega(h)$.
These are standard convexity estimates, also used in the local
comparison framework of~\cite{BEY}. One way to verify the logarithmic
Sobolev constant is to use the reflecting diffusion
with generator $\Delta-\nabla\Phi\cdot\nabla$ on an inner polytope
after smooth convex approximation. The Bochner inequality, including
its nonnegative boundary term, gives decay of Fisher information at
rate $2/T$; integrating the entropy dissipation identity yields
$S\le(T/2)I$.
The passage to the full chamber does not require a boundary
integrability threshold on $\beta$. To see this directly, first take
a bounded locally Sobolev function $f$ with finite weighted Dirichlet
integral. Its restriction to each inner polytope belongs to the usual
Sobolev space, because the weight is smooth and bounded above and
below there. Apply its logarithmic Sobolev inequality and let
$\varepsilon\downarrow0$. The entropy integrals converge by
boundedness, and the Dirichlet integrals converge by monotone
convergence before normalization. Truncation of an arbitrary
finite-energy $f$, followed by lower semicontinuity of entropy, proves
the general assertion. Apply this to
$f=\sqrt{\dd\mu/\dd\omega_{K,T}}$ and combine with
\eqref{sta:entropy-gap}. All preceding constants depend on $\beta$
only through positive lower bounds for $\beta m$, proving the stated
uniformity.
\end{proof}

\section{A temperature comparison for convex reference measures}
\label{sta:sec-temperature}

The homogeneity of an inverse-power interaction changes its effective
inverse temperature under dilation. The following finite-dimensional
estimate controls that change without estimating singular pair forces.
Its variance bound is the sharp information-content variance inequality
for log-concave densities; we follow the perspective proof presented
in~\cite{FMW}. The argument is included for completeness.

\begin{lemma}[Convex temperature bound]\label{sta:temperature}
Let $D\subset\R^K$ be a convex set with nonempty interior, and let
$\Phi$ be a measurable convex function on $D$, finite almost everywhere
there and bounded below. Extend it by $+\infty$ off $D$. Assume
\[
0<Z(t):=\int_{\R^K}e^{-t\Phi(u)}\dd u<\infty
\qquad(t>0),
\]
and put $\omega_t(\dd u)=Z(t)^{-1}e^{-t\Phi(u)}\dd u$.
Then, for every $t>0$ and $\tau>0$,
\begin{align}
\Var_{\omega_t}(\Phi)&\le K/t^2,
 \label{sta:energy-variance}\\
S(\omega_\tau\mid\omega_1)
 &\le K(\log\tau+\tau^{-1}-1).
 \label{sta:temperature-entropy}
\end{align}
\end{lemma}

\begin{proof}
Adding a constant to $\Phi$ changes none of the measures or claimed
quantities, so assume $\Phi\ge0$. For $t>0$ the perspective
\[
(y,t)\longmapsto t\Phi(y/t)
\]
is convex as an extended-valued function. Indeed, if
$t=(1-\theta)t_0+\theta t_1$, then
$y/t$ is the convex combination of $y_0/t_0,y_1/t_1$ with weights
$(1-\theta)t_0/t$ and $\theta t_1/t$, respectively. Multiplication of
the convexity inequality for $\Phi$ by $t$ proves the claim, including
the constraint $y/t\in D$.
The Pr\'ekopa marginal log-concavity theorem implies that
\[
t\longmapsto\log\int_{\R^K}e^{-t\Phi(y/t)}\dd y
 =K\log t+\log Z(t)
\]
is concave. Let $\psi(t)=\log Z(t)$. Differentiating under the integral
is justified at every $t>0$: for $j=1,2$ and $t$ in a small compact
neighbourhood of $t_0>0$, the factors $\Phi^j e^{-t\Phi}$ are dominated
by a constant times $e^{-(t_0/2)\Phi}$ after shrinking that
neighbourhood. Consequently
\[
\psi'(t)=-\omega_t(\Phi),\qquad
\psi''(t)=\Var_{\omega_t}(\Phi).
\]
The second derivative of the concave function above is nonpositive,
and gives \eqref{sta:energy-variance}.

The logarithmic density ratio gives the exact identity
\begin{equation}\label{sta:kl-identity}
S(\omega_\tau\mid\omega_1)
=\psi(1)-\psi(\tau)+(\tau-1)\psi'(\tau).
\end{equation}
Its derivative in $\tau$ is $(\tau-1)\psi''(\tau)$, and its value at
$\tau=1$ is zero. Thus, when $\tau\ge1$,
\[
S(\omega_\tau\mid\omega_1)
=\int_1^\tau(t-1)\psi''(t)\dd t
\le K\int_1^\tau\frac{t-1}{t^2}\dd t
=K(\log\tau+\tau^{-1}-1).
\]
When $0<\tau\le1$, the positive-kernel expression is instead
\[
S(\omega_\tau\mid\omega_1)
=\int_\tau^1(1-t)\psi''(t)\dd t
\le K\int_\tau^1\frac{1-t}{t^2}\dd t
=K(\log\tau+\tau^{-1}-1).
\]
The change of integration orientation is essential for using the
variance upper bound in this second case.
\end{proof}

\begin{corollary}[Temperature bridge for gap averages]
\label{sta:temperature-gap}
Let $\Phi=\Phi_{K,T}$ be as in \eqref{sta:reference} and define
$\omega_{K,T}^{(\tau)}\propto e^{-\tau\Phi_{K,T}}$.
Uniformly in $K\ge2r+1$ and $T>0$,
\begin{equation}\label{sta:temperature-observable}
|\omega_{K,T}^{(\tau)}(F_K)-\omega_{K,T}(F_K)|^2
\le C(\log\tau+\tau^{-1}-1).
\end{equation}
In particular, on each fixed compact interval
$J\subset(0,\infty)$ containing $1$, its right-hand side is at most
$C_J|\tau-1|^2$.
\end{corollary}

\begin{proof}
Convexity on the bounded chamber gives a finite lower bound for
$\Phi_{K,T}$, so all hypotheses of Lemma~\ref{sta:temperature} hold.
Use \eqref{sta:entropy-gap} with
$\mu=\omega_{K,T}^{(\tau)}$, and then
\eqref{sta:temperature-entropy}. The factor $K$ cancels. For the final
statement, write $h(\tau)=\log\tau+\tau^{-1}-1$; one has
$h(1)=h'(1)=0$ and $h''$ bounded on $J$, so Taylor's formula applies.
\end{proof}

\begin{remark}\label{sta:whole-potential}
The temperature parameter multiplies the entire potential:
\[
\tau\Phi_{K,T}
=\tau\beta H^{\mathrm{lat}}_{V,K}
  +\frac{1}{2(T/\tau)}\sum_i(u_i-i)^2.
\]
Thus the bridge changes the interaction coefficient to $\tau\beta$
and the confinement parameter to $T/\tau$. This is the family for which
Lemma~\ref{sta:temperature} applies directly. Multiplying just one
part of the Hamiltonian would require a separate argument.
\end{remark}


\section{Stationary counting and particle-ended blocks}
\label{pre:section}\label{sec:prelim}

This section uses no special property of the interaction. Let $P$ be a
stationary simple point process of intensity one and write
$v_P(t)=\operatorname{Var}_P N_{(0,t]}$. Assume throughout that
\begin{equation}\label{pre:variance}
v_P(t)\le C_Pt\quad(t>0),\qquad v_P(t)=o(t)\quad(t\to\infty).
\end{equation}
Translations are $\theta_x\gamma=\gamma-x$. The Palm law $P^0$ is
characterized by Campbell's identity
\begin{equation}\label{pre:campbell}
\E_P\sum_{p\in\gamma} h(p,\theta_p\gamma)
=\int_{\R}\E_{P^0}h(x,\gamma)\dd x
\end{equation}
for nonnegative measurable $h$. Under $P^0$ enumerate the points as
$\cdots<p_{-1}<p_0=0<p_1<\cdots$, put $g_i=p_{i+1}-p_i$, and let
$T_*\gamma=\theta_{p_1}\gamma$.

\begin{lemma}[Density one]\label{pre:density}
Almost surely under both $P$ and $P^0$, the configuration has density
one in both directions, at every finite origin.
\end{lemma}

\begin{proof}
Chebyshev's inequality gives
$P(|N_{(0,n^2]}-n^2|>\varepsilon n^2)\le C_P/(\varepsilon^2n^2)$.
Borel--Cantelli and monotonic interpolation prove forward density one;
the reflected argument proves backward density one. Moving the origin
by a finite amount does not change these limits. Thus the good event
is invariant under every translation, and Campbell's identity
transfers it to $P^0$.
\end{proof}

\begin{lemma}[Palm count envelopes]\label{pre:envelopes}
The random variables
\[
C^+=\sup_{t\ge0}\frac{N_{(0,t]}}{1+t},\qquad
C^-=\sup_{t\ge0}\frac{N_{[-t,0)}}{1+t}
\]
are integrable under $P^0$.
\end{lemma}

\begin{proof}
For a spatially rooted configuration put
$C_0=\sup_{t\ge0}N_{(0,t]}/(1+t)$. A dyadic decomposition gives
\[
C_0\le N_{(0,1]}+2+
\sum_{j\ge0}2^{-j}|N_{(0,2^{j+1}]}-2^{j+1}|.
\]
Minkowski's inequality and the linear variance bound imply
$\E_PC_0^2<\infty$. For every $p\in\gamma\cap(0,1]$,
$C^+(\theta_p\gamma)\le2C_0(\gamma)$. Consequently, by Campbell and
Cauchy--Schwarz,
\[
\E_{P^0}C^+
=\E_P\sum_{p\in\gamma\cap(0,1]}C^+(\theta_p\gamma)
\le2\E_P[N_{(0,1]}C_0]<\infty.
\]
Reflection proves the assertion for $C^-$.
\end{proof}

\begin{lemma}[First-point bias and recovery of the Palm law]
\label{pre:palm-identification}
The Palm law is invariant under $T_*$. If $a$ is the first strictly
positive point of a spatial sample, then
\begin{equation}\label{pre:first-point}
\operatorname{Law}_P(\theta_a\gamma)
=g_{-1}P^0,\qquad \E_{P^0}g_{-1}=1.
\end{equation}
For every bounded measurable function $G$ of $r$ consecutive gaps,
the expectation of its average over successive gap windows following
$a$ converges, as the number of windows tends to infinity, to
$\E_{P^0}G(g_0,\ldots,g_{r-1})$.
Moreover the Palm law determines $P$ through
\begin{equation}\label{pre:inversion}
\E_P f(\gamma)
=\E_{P^0}\int_0^{g_0}f(\theta_t\gamma)\dd t
\end{equation}
for bounded measurable $f$.
\end{lemma}

\begin{proof}
For bounded $h$, compare the sums of $h(\theta_p\gamma)$ and
$h(T_*\theta_p\gamma)$ over $p\in\gamma\cap(0,R]$. Their sets of
point indices differ at most at the two ends, so the absolute
difference is at most $2\norm h_\infty$. Campbell and division by $R$
give $\E_{P^0}h(T_*\gamma)=\E_{P^0}h(\gamma)$.

There is a unique point immediately to the left of the spatial
origin, by Lemma~\ref{pre:density}; a point at zero has $P$-probability
zero. Assigning the origin to that point and applying Campbell gives
\[
\E_P f(\gamma)
=\int_{\R}\E_{P^0}
 [\1_{\{x\le0<x+g_0\}}f(\theta_{-x}\gamma)]\dd x
=\E_{P^0}\int_0^{g_0}f(\theta_t\gamma)\dd t.
\]
Assigning the origin instead to the first point on its right gives
\[
\E_P h(\theta_a\gamma)
=\int_{\R}\E_{P^0}
 [\1_{\{0<x\le g_{-1}\}}h(\gamma)]\dd x
=\E_{P^0}[g_{-1}h].
\]
Taking $f=1$ or $h=1$ gives the mean gap identity.

For completeness the averaging assertion does not require ergodicity.
Let $\mathcal I$ be the $T_*$-invariant sigma-algebra. Birkhoff's
theorem gives the almost-sure limit
$\E_{P^0}[G\mid\mathcal I]$ for the bounded gap averages.
The density assertion and Birkhoff applied to $g_{-1}$ give
\[
\E_{P^0}[g_{-1}\mid\mathcal I]
=\lim_{n\to\infty}\frac1n\sum_{j=0}^{n-1}g_{-1}\circ T_*^j=1.
\]
Use \eqref{pre:first-point}, dominated convergence with the integrable
weight $g_{-1}$, and conditional expectation onto $\mathcal I$.
\end{proof}

\begin{remark}[Endpoint domination]\label{pre:domination}
Let $g_*$ denote the entire gap containing the spatial origin.
For every nonnegative measurable $h$, every integer $j$, and $M>0$,
\begin{equation}\label{pre:domination-equation}
\E_P[\1_{\{g_*\le M\}}h(T_*^j\theta_a\gamma)]
=\E_{P^0}[g_{-1}\1_{\{g_{-1}\le M\}}h(T_*^j\gamma)]
\le M\E_{P^0}h.
\end{equation}
The bound is uniform in $j$. In particular it applies to the two
particle endpoints of all blocks considered below.
\end{remark}

\subsection{Canonical conditioning on a fixed-rank block}

For $K\ge1$ let $a$ be the first positive point and $b$ the
$(K+2)$-nd positive point. Write the $K$ points between them as
$x_1<\cdots<x_K$ and define
\[
L=K+1,\qquad \ell=b-a,\qquad c=L/\ell,\qquad
u_i=c(x_i-a),\qquad
\zeta=\gamma-\sum_{i=1}^K\delta_{x_i}.
\]
Thus $a,b$ are the first two positive points of the retained
configuration $\zeta$, and $u\in\Omega_K=\{0<u_1<\cdots<u_K<L\}$.

\begin{lemma}[Random-rank DLR disintegration]\label{pre:random-dlr}
Suppose $P$ satisfies the canonical DLR equations for the interaction
$V$. Conditionally on $\zeta$, the points $x_1,\ldots,x_K$ have the
canonical Gibbs density in $(a,b)$ with exterior $\zeta$. In the
$u$-coordinates the conditional law, denoted by $\mu_\zeta$, has
interaction $V_c(r)=V(r/c)$ and exterior $c(\zeta-a)$, including
charges at $0$ and $L$.
\end{lemma}

\begin{proof}
Restrict first to $\{b<R\}$, where $R$ is a positive integer. Apply
canonical DLR in the deterministic interval $(0,R)$, conditional on
its exterior and particle number. Then condition further on the
ordered coordinates of rank $1,K+2,K+3,\ldots$ in that interval.
Only the $K$ coordinates between $a$ and $b$ remain variable.
Factoring their terms out of the original Gibbs density gives exactly
their Gibbs density with the retained exterior. Both the total number
in $(0,R)$ and the event $\{b<R\}$ are determined by $\zeta$ and $K$.
The rank selection is preserved by every resampling inside $(a,b)$,
so it introduces no additional coordinate-dependent factor.
The events $\{b<R\}$ exhaust probability one. Finally the affine
change of variables contributes a constant Jacobian and transforms
each distance $r$ into $r/c$. For the relative exterior energies used
here the same change is legitimate by absolute tail convergence for
inverse powers and by the symmetric-cutoff convergence for the
logarithm. A finite displacement of the cutoff center contributes
vanishing distant-shell terms by density one.
\end{proof}

\begin{lemma}[Good lengths]\label{pre:good-length}
One has $\ell-L=o_P(\sqrt K)$. For any two candidate laws $P,Q$,
there is a common deterministic sequence $r_K\downarrow0$ such that,
for every fixed $M>0$, the exterior-measurable event
\begin{equation}\label{pre:good-event}
E_{K,M}=\{g_*\le M,\ |\ell-L|\le r_K\sqrt K\}
\end{equation}
satisfies
\[
\limsup_{K\to\infty}P(E_{K,M}^c)\le P(g_*>M),
\]
and the analogous inequality holds under $Q$.
\end{lemma}

\begin{proof}
Put $n=K+2$. For $t=n+\varepsilon\sqrt K$, the event $\{b>t\}$
implies $N_{(0,t]}<n$, a deviation of size at least
$\varepsilon\sqrt K$. Since $v_P(t)=o(K)$, its probability tends to
zero by Chebyshev. Taking $t=n-\varepsilon\sqrt K$ similarly controls
$\{b<t\}$. Therefore $b-n=o_P(\sqrt K)$. The random variable $a$
is almost surely finite and has a law independent of $K$, proving the
claim about $\ell-L=b-n+1-a$. A diagonal choice gives a common slowly
decreasing $r_K$ for $P,Q$. The other assertions follow because $a,b$
and the last point before zero all belong to $\zeta$.
\end{proof}

\begin{lemma}[Negligible confinement cost]\label{pre:quadratic}
Define
\begin{equation}\label{pre:J-definition}
J_K(\zeta)=\E_{\mu_\zeta}\sum_{i=1}^K(u_i-i)^2.
\end{equation}
For every fixed $M$,
\begin{equation}\label{pre:J-small}
\frac1{K^2}\E_P[\1_{E_{K,M}}J_K]\longrightarrow0.
\end{equation}
\end{lemma}

\begin{proof}
For any deterministic $0<z_1<\cdots<z_K<\ell$, let
$A(t)=\#\{i:z_i\le t\}$ and $\rho=K/\ell$. The exact identity
\begin{equation}\label{pre:quantile}
\int_0^\ell(A(t)-\rho t)^2\dd t
=\rho\sum_{i=1}^K
\left(z_i-\frac{i-\frac12}{\rho}\right)^2+\frac\ell{12}
\end{equation}
follows on expanding the square and using
\[
\int_0^\ell A(t)^2\dd t=\sum_i(2i-1)(\ell-z_i),\qquad
\int_0^\ell tA(t)\dd t=\frac12\sum_i(\ell^2-z_i^2).
\]
Take $z_i=x_i-a$ and $D_0(s)=N_{(0,s]}-s$. For Lebesgue-almost
every $0<t<\ell$,
\[
A(t)-\rho t=D_0(a+t)+(a-1)+(1-\rho)t.
\]
On $E_{K,M}$ we have $a\le g_*\le M$, $\ell\asymp K$, and $b\le3K$
for sufficiently large $K$. The midpoint locations
$(i-\frac12)\ell/K$ differ from $i\ell/L$ by a uniformly bounded
amount. Therefore \eqref{pre:quantile} yields
\begin{equation}\label{pre:quadratic-pointwise}
\1_{E_{K,M}}\sum_i(u_i-i)^2
\le C\int_0^{3K}D_0(s)^2\dd s+C_MK+Cr_K^2K^2.
\end{equation}
Here we used $K(\ell-K)^2\le C K+C K(\ell-L)^2$ for the density
correction. By stationarity and intensity one,
$\E_PD_0(s)^2=v_P(s)$, and
$\int_0^{3K}v_P(s)\dd s=o(K^2)$. Take expectations in
\eqref{pre:quadratic-pointwise}. Exterior measurability of $E_{K,M}$
allows conditional expectation of the left side, giving
\eqref{pre:J-small}.
\end{proof}

\begin{lemma}[Conditional envelopes]\label{pre:conditional-envelopes}
For an interior configuration define
\[
C_{\rm int}^+=\sup_{t\ge0}\frac{\#\{i:u_i\le t\}}{1+t},\qquad
C_{\rm int}^-=\sup_{t\ge0}\frac{\#\{i:L-u_i\le t\}}{1+t},\qquad
X_\pm=\E_{\mu_\zeta}C_{\rm int}^\pm.
\]
For each fixed $M$ and all sufficiently large $K$,
\begin{equation}\label{pre:X-bound}
\E_P[\1_{E_{K,M}}X_\pm]\le2M\E_{P^0}C^\pm.
\end{equation}
Furthermore, for any nonnegative decreasing function $h$ on $[0,L]$,
\begin{equation}\label{pre:decreasing-envelope}
\sum_{i=1}^K h(u_i)^2
\le C_{\rm int}^+\left(h(0)^2+\int_0^Lh(t)^2\dd t\right),
\end{equation}
and the reflected inequality holds at the right endpoint.
\end{lemma}

\begin{proof}
On $E_{K,M}$, $c\in[1/2,2]$, and changing variables in the count
envelope gives
$C_{\rm int}^+\le2C^+(\theta_a\gamma)$ and
$C_{\rm int}^-\le2C^-(\theta_b\gamma)$.
Condition on $\zeta$, apply \eqref{pre:domination-equation} with
$j=0,L$, and use Lemma~\ref{pre:envelopes}.
For \eqref{pre:decreasing-envelope}, integrate $h^2$ against the
interior counting function $N(t)$ and use
$N(t)\le C_{\rm int}^+(1+t)$. Integration by parts gives
\[
\int_0^Lh^2\dd N
\le C_{\rm int}^+\left((1+L)h(L)^2+
\int_{[0,L]}(1+t)\,\dd(-h(t)^2)\right)
=C_{\rm int}^+\left(h(0)^2+\int_0^Lh^2\dd t\right).
\]
Approximation proves the assertion for arbitrary decreasing $h$.
\end{proof}

\begin{lemma}[Affine block statistics]\label{pre:affine-statistics}
Let $G\in C_c^2((0,\infty)^r)$ and
\[
F_K(u)=\frac1{K-r}\sum_{i=1}^{K-r}
G(u_{i+1}-u_i,\ldots,u_{i+r}-u_{i+r-1}).
\]
Then
\begin{equation}\label{pre:affine-limit}
\lim_{K\to\infty}\E_P F_K(u(\gamma))
=\E_{P^0}G(g_0,\ldots,g_{r-1}).
\end{equation}
Consequently equality of these limits for two intensity-one laws and
all $r,G$ implies equality of the laws.
\end{lemma}

\begin{proof}
Extend $G$ by zero to a globally Lipschitz function on $\R^r$.
Each physical gap enters at most $r$ summands, and the sum of the
interior gaps is at most $\ell$. Thus the difference between the
scaled statistic and its unscaled version is at most
\[
\frac{C_G}{K-r}|c-1|\ell
=\frac{C_G}{K-r}|\ell-L|\longrightarrow0
\]
almost surely, by density one. Both statistics are bounded, so the
difference tends to zero in mean. Lemma~\ref{pre:palm-identification}
identifies the unscaled limit. Equality for every $G$ identifies all
finite consecutive-gap distributions; $T_*$-invariance identifies
every finite window of integer gap indices. The full gap sequence
determines the Palm configuration, and \eqref{pre:inversion} recovers
the spatial law.
\end{proof}


\section{Exterior forces and tight conditional costs}\label{for:section}\label{sec:forces}

The following argument treats the logarithm and the inverse powers
together.  Fix $\alpha\ge0$ and write
\[
 V_\alpha(r)=
 \begin{cases}-\log r,&\alpha=0,\\ r^{-\alpha},&\alpha>0.
 \end{cases}
\]
All constants in this section may depend on $\alpha$.  Let $P$ be a
stationary simple point process of intensity one, and put
$v(t)=\operatorname{Var}_P N_{(0,t]}$.  We use the assumptions
\begin{equation}\label{for:assumptions}
 v(t)\le Ct\quad(t>0),\qquad v(t)=o(t),\qquad
 \int_1^\infty v(t)t^{-2\alpha-2}\,\dd t<\infty.
\end{equation}
For every $\alpha>0$ the last condition follows from the first, since
$\int_1^\infty t^{-2\alpha-1}\,\dd t<\infty$.  For $\alpha=0$ it
is the additional logarithmic discrepancy assumption.  The sublinear
condition is only needed below through the good endpoint-length event.

\begin{lemma}[Weighted discrepancy regularity]\label{for:regularity}
There is a translation-invariant event of full $P$- and $P^0$-measure
on which the configuration has density one in both directions and,
simultaneously for every $a\in\R$,
\begin{equation}\label{for:weighted-discrepancy}
 \int_1^\infty
 \bigl((N_{(a,a+t]}-t)^2+(N_{[a-t,a)}-t)^2\bigr)
 t^{-2\alpha-2}\,\dd t<\infty.
\end{equation}
\end{lemma}

\begin{proof}
Density one was proved in Lemma~\ref{pre:density}.
Tonelli's theorem and \eqref{for:assumptions} give
\eqref{for:weighted-discrepancy} at $a=0$.  To pass simultaneously to
all finite $a$, use the signed cumulative discrepancy $D_0$ of
$\gamma-\dd x$.  An interval discrepancy rooted at $a$ is a difference
$D_0(a+t)-D_0(a)$, up to a fixed endpoint convention.  For large $t$,
$t$ and $t+a$ are comparable.  Translation in the integral and
$\int_1^\infty t^{-2\alpha-2}\,\dd t<\infty$ prove the assertion
at every $a$ on the same event.  Density one is likewise unaffected by
finite translations.  This event is consequently translation invariant;
Campbell's formula transfers its probability-one property from $P$ to
$P^0$.
\end{proof}

\begin{lemma}[Logarithmic principal values]\label{for:log-pv}
When $\alpha=0$, the event in Lemma~\ref{for:regularity} also has
the following property. At every finite origin the reciprocal
symmetric principal value exists, with a point at that origin omitted.
The logarithmic relative exterior energies converge, and their first
and second derivatives converge locally uniformly away from exterior
points. A finite change of cutoff center does not change their limits.
\end{lemma}

\begin{proof}
At origin zero put
$D_{\rm sym}(t)=N_{(0,t]}-N_{[-t,0)}$. Stieltjes integration by parts gives
\[
\sum_{\substack{p\in\gamma\\1<|p|\le R}}\frac1p
=\frac{D_{\rm sym}(R)}R-D_{\rm sym}(1)
+\int_1^R\frac{D_{\rm sym}(t)}{t^2}\dd t.
\]
Density one makes the boundary term tend to zero.
The integral converges absolutely by Cauchy--Schwarz and
\eqref{for:weighted-discrepancy}. The same argument applies at every
finite origin on the same event.
For $x,x_*$ in a fixed bounded set and large $|p|$,
\[
-\log|x-p|+\log|x_*-p|
=\frac{x-x_*}{p}+O(|p|^{-2}),
\]
uniformly on that set. The reciprocal principal value and the
absolutely summable remainder prove the relative-energy claim.
The derivative has the same principal-value leading term, and the
second derivative is absolutely summable.
Finally, moving a cutoff center by a fixed amount changes only shells
of bounded width. Density one bounds the number of points in those
shells by $o(R)$, and each relative-energy or first-derivative summand
is $O(R^{-1})$. Their contribution tends to zero.
\end{proof}

\begin{lemma}[Half-line envelope and its dilation]\label{for:halfline}
Root a configuration from the preceding event at a particle.  Remove
that particle, and let $N_{\rm ext}(t)$ count the particles on either
chosen outward half-line at distance at most $t$.  Define
\begin{equation}\label{for:envelope-def}
 D(t)=N_{\rm ext}(t)-(t-1)_+,\qquad
 H_\alpha(x)=\int_0^\infty
 \frac{\abs{D(t)}}{(x+t)^{\alpha+2}}\,\dd t\quad(x\ge0).
\end{equation}
Then
\begin{equation}\label{for:envelope-energy}
 \mathcal B_\alpha
 :=H_\alpha(0)^2+\int_0^\infty H_\alpha(x)^2\,\dd x<\infty.
\end{equation}
More precisely, with $\mathrm B$ denoting the beta function,
\begin{equation}\label{for:schur}
 \int_0^\infty H_\alpha(x)^2\,\dd x
 \le \mathrm B\left(\frac12,\alpha+\frac32\right)^2
 \int_0^\infty D(t)^2t^{-2\alpha-2}\,\dd t.
\end{equation}
If outward distances are multiplied by $c\in[1/2,2]$, center the new
counting function at $c^{-1}(t-1)_+$.  The resulting envelope energy
$\mathcal B_{\alpha,c}$ satisfies
\begin{equation}\label{for:dilation}
 \sup_{1/2\le c\le2}\mathcal B_{\alpha,c}
 \le C_\alpha(\mathcal B_\alpha+1).
\end{equation}
\end{lemma}

\begin{proof}
Simplicity and local finiteness give a positive distance to the next
outward point.  Hence $D$ vanishes on some interval $(0,\delta)$.
At infinity $D(t)$ differs from $N_{\rm ext}(t)-t$ by the constant
$1$, so Lemma~\ref{for:regularity} gives
$\int_0^\infty D(t)^2t^{-2\alpha-2}\,\dd t<\infty$.
Cauchy--Schwarz yields
\[
 H_\alpha(0)
 \le\left(\int_\delta^\infty D(t)^2t^{-2\alpha-2}\,\dd t\right)^{1/2}
     \left(\int_\delta^\infty t^{-2}\,\dd t\right)^{1/2}<\infty.
\]
For the square-integral bound, write
$f(t)=\abs{D(t)}t^{-\alpha-1}$ and use the positive kernel
\[
 k_\alpha(x,t)=\frac{t^{\alpha+1}}{(x+t)^{\alpha+2}}.
\]
With $w(t)=t^{-1/2}$, direct substitutions give both Schur identities
\[
 \int_0^\infty k_\alpha(x,t)w(t)\,\dd t
 =b_\alpha w(x),\qquad
 \int_0^\infty k_\alpha(x,t)w(x)\,\dd x
 =b_\alpha w(t),
 \quad b_\alpha=\mathrm B\left(\tfrac12,\alpha+\tfrac32\right).
\]
The Schur bound proves \eqref{for:schur}.  In particular
$b_0^2=\pi^2/4$.

For the dilation, the new centered discrepancy is
\[
 D_c(t)=D(t/c)+R_c(t),\qquad
 R_c(t)=(t/c-1)_+-c^{-1}(t-1)_+.
\]
Uniformly for $c\in[1/2,2]$, $R_c$ is bounded and vanishes when
$t<1/2$.  Therefore
\[
 H_{\alpha,c}(x)
 \le c^{-\alpha-1}H_\alpha(x/c)
       +C_\alpha(x+1/2)^{-\alpha-1}.
\]
Its value at zero and its square integral are bounded as claimed in
\eqref{for:dilation}.
\end{proof}

Retain the random block notation: $a$ is the first positive point,
$b$ is the $(K+2)$-nd positive point, $L=K+1$, $\ell=b-a$,
$c=L/\ell$, and $u_i=c(x_i-a)$.  The retained exterior $\zeta$
includes $a,b$, and $\mu_\zeta$ denotes the conditional distribution
of the rescaled interior.  Write $\widehat\zeta=c(\zeta-a)$ and
\[
 \xi_K^{\rm lat}=\{0,-1,-2,\ldots\}\cup\{L,L+1,L+2,\ldots\}.
\]
The common endpoint charges $0,L$ will always be cancelled before
estimating a force difference.  To fix a harmless sign convention, put
$q_\alpha(r)=-V_\alpha'(r)$ and define
\begin{equation}\label{for:force-def}
 \begin{split}
 B_\zeta(u)={}&
 \sum_{p\in\widehat\zeta\setminus\{0,L\}}
     \operatorname{sgn}(u-p)q_\alpha(\abs{u-p})\\
 &-\sum_{p\in\xi_K^{\rm lat}\setminus\{0,L\}}
     \operatorname{sgn}(u-p)q_\alpha(\abs{u-p}),\qquad 0\le u\le L.
 \end{split}
\end{equation}
For $\alpha>0$ these force sums are absolutely convergent.  For
$\alpha=0$ each sum is understood in symmetric principal value, or
equivalently by pairing the two outward rays.  Indeed, subtraction of
the corresponding constant-density rays leaves absolutely convergent
discrepancy integrals, and the paired continuum force converges.
Changing between the two cutoff conventions contributes vanishing
distant shells by density one.  The estimates below in particular
show that $B_\zeta$ is bounded on the closed block for almost every
retained exterior.

Define its conditional cost by
\begin{equation}\label{for:cost}
 \mathcal D_{K,\alpha}(\zeta)=\mathcal D_K(\zeta)
 =\E_{\mu_\zeta}\sum_{i=1}^K B_\zeta(u_i)^2.
\end{equation}

\begin{proposition}[Tight conditional force cost]\label{for:tight}
Let $r_K\downarrow0$ be chosen so that
$\ell-L=o_P(\sqrt K)$ is captured by the good events
\[
 E_{K,M}=\{g_*\le M,\ \abs{\ell-L}\le r_K\sqrt K\},
\]
where $g_*$ is the full gap containing the spatial origin, as in
Lemma~\ref{pre:good-length}.
For each fixed $M$, the random variables
$\1_{E_{K,M}}\mathcal D_K$ are tight as $K\to\infty$.
\end{proposition}

\begin{proof}
On $E_{K,M}$, $c\in[1/2,2]$ for all sufficiently large $K$.
The rescaled exterior has density $\lambda=c^{-1}=\ell/L$.
For each of its two outward rays, center its residual count at
$\lambda(t-1)_+$ and denote the resulting envelopes and energies
by $H_{\alpha,+},H_{\alpha,-}$ and
$\mathcal B_{\alpha,+},\mathcal B_{\alpha,-}$.
These are functions of the retained exterior.

Stieltjes integration by parts compares a ray force to its continuum
counterpart.  Since
$\abs{q_\alpha'(r)}=C_\alpha r^{-\alpha-2}$, its discrepancy
contribution is bounded by $C_\alpha H_{\alpha,+}(u)$, or by its
reflection at the right endpoint.  The boundary term at infinity
vanishes: the centered count is $o(t)$, and
$q_\alpha(u+t)=O(t^{-\alpha-1})$.  For the unit lattice, the residual
count is $\lfloor t\rfloor$.  Its difference from $(t-1)_+$ is
bounded and vanishes below $1$, yielding the deterministic envelope
$C_\alpha(u+1)^{-\alpha-1}$.

The continuum density mismatch is explicit.  Up to the overall sign
chosen in \eqref{for:force-def}, its contribution is
$(\lambda-1)f_{\alpha,L}(u)$, where
\[
 f_{\alpha,L}(u)=
 \begin{cases}
 \displaystyle\log\frac{L-u+1}{u+1},&\alpha=0,\\[6pt]
 (u+1)^{-\alpha}-(L-u+1)^{-\alpha},&\alpha>0.
 \end{cases}
\]
Consequently
\begin{equation}\label{for:pointwise}
 \begin{split}
 \abs{B_\zeta(u)}\le C_\alpha\bigl\{&
 H_{\alpha,+}(u)+H_{\alpha,-}(L-u)
 +(u+1)^{-\alpha-1}+(L-u+1)^{-\alpha-1}\bigr\}\\
 &+\abs{\lambda-1}\abs{f_{\alpha,L}(u)}.
 \end{split}
\end{equation}
There are no endpoint singularities in this estimate, because the
charges at $0,L$ have already cancelled exactly.

Recall $C_{\rm int}^\pm$ and $X_\pm$ from
Lemma~\ref{pre:conditional-envelopes}.  Its decreasing-function
estimate gives
\begin{equation}\label{for:summation}
 \sum_i h(u_i)^2
 \le C_{\rm int}^+
 \left(h(0)^2+\int_0^L h(t)^2\,\dd t\right),
\end{equation}
and an identical reflected estimate.

For $\alpha=0$ take
$h_{0,L}(t)=\log((L+1)/(t+1))$, and for $\alpha>0$ take
$h_{\alpha,L}(t)=(t+1)^{-\alpha}$.  In either case
$\abs{f_{\alpha,L}(u)}\le h_{\alpha,L}(u)+h_{\alpha,L}(L-u)$.
Moreover, for $L\ge2$,
\[
 h_{\alpha,L}(0)^2+\int_0^L h_{\alpha,L}(t)^2\,\dd t
 \le C_\alpha A_\alpha(L),\qquad
 A_\alpha(L)=
 \begin{cases}
 L,&\alpha=0,\\
 \displaystyle1+\int_0^L(1+t)^{-2\alpha}\,\dd t,&\alpha>0.
 \end{cases}
\]
For the logarithm, the integral is at most
$(L+1)\int_0^1\log(1/y)^2\,\dd y=2(L+1)$ and
$\log(L+1)^2\le CL$.  For positive $\alpha$ this assertion is
immediate; its growth is $L^{1-2\alpha}$, $\log L$, or $1$
according as $\alpha<1/2$, $\alpha=1/2$, or $\alpha>1/2$.
Squaring \eqref{for:pointwise}, applying \eqref{for:summation}, and
taking the conditional expectation now gives
\begin{equation}\label{for:cost-bound}
 \begin{split}
 \mathcal D_K\le C_\alpha\bigl\{&
 X_+(\mathcal B_{\alpha,+}+1)
 +X_-(\mathcal B_{\alpha,-}+1)\\
 &+(\lambda-1)^2A_\alpha(L)(X_++X_-)\bigr\}.
 \end{split}
\end{equation}

We finally check tightness under the random endpoint law.
The endpoint domination \eqref{pre:domination-equation} bounds the
laws rooted at either $a$ or $b$, restricted to $\{g_*\le M\}$,
by $M P^0$, uniformly in $K$.
Lemma~\ref{for:halfline}, including its uniform dilation bound,
therefore implies that the two exterior energies restricted to
$E_{K,M}$ are tight.  The restricted $X_\pm$ are tight by
\eqref{pre:X-bound}; that estimate uses integrability of the full
Palm count envelopes and exterior measurability of $E_{K,M}$.
Products of two tight
nonnegative families are tight: the event that a product exceeds
$R^2$ is contained in the union of the events that either factor
exceeds $R$.  No independence or integrable product is required.
Finally $A_\alpha(L)\le C_\alpha L$ and, on $E_{K,M}$,
\[
 (\lambda-1)^2A_\alpha(L)
 \le C_\alpha\frac{(\ell-L)^2}{L}
 \le C_\alpha r_K^2\longrightarrow0.
\]
Insert these conclusions into \eqref{for:cost-bound} to obtain the
proposition.
\end{proof}


\section{Completion of the comparison arguments}\label{sec:completion}

\subsection{A common-reference criterion}
Fix two candidate laws $P,Q$, an integer $r\ge1$ and
$G\in C_c^2((0,\infty)^r)$, and use the random blocks and good events
of Section~\ref{sec:prelim}. For each fixed $M$, choose a common
sequence $r_K\downarrow0$ for the two processes.

\begin{lemma}[Closing the comparison]\label{end:criterion}
Suppose a deterministic reference $\omega_{K,T}$ satisfies, on
$E_{K,M}$ and for $K$ sufficiently large,
\begin{equation}\label{end:abstract}
\abs{\mu_\zeta(F_K)-\omega_{K,T}(F_K)}^2
\le C_G\left\{\frac{T}{K}\mathcal A_K(\zeta)
+\frac{J_K(\zeta)}{KT}+e_K(\zeta)^2\right\},
\end{equation}
where $\mathcal A_K,J_K\ge0$ and, under each candidate law, $\1_{E_{K,M}}\mathcal A_K$ is tight,
$\E[\1_{E_{K,M}}J_K]=o(K^2)$, and
$\sup_{E_{K,M}}|e_K|\to0$. The constants are independent of $K,T$.
If these estimates hold for every $r,G,M$, then $P=Q$.
\end{lemma}

\begin{proof}
For fixed $M$, put
\[
q_K=\frac{\E_P[\1_{E_{K,M}}J_K]
+\E_Q[\1_{E_{K,M}}J_K]}{K^2},\qquad
\varepsilon_K=\max\{\sqrt{q_K},K^{-1/2}\},\qquad
T_K=\varepsilon_KK.
\]
Both $q_K$ and $\varepsilon_K$ tend to zero, and
$q_K/\varepsilon_K\le\sqrt{q_K}$. At this choice of $T$ the
right-hand side of \eqref{end:abstract} is
\[
C_G\left\{\varepsilon_K\mathcal A_K+
\frac{J_K}{\varepsilon_KK^2}+e_K^2\right\}.
\]
Restricted to $E_{K,M}$, its first term tends to zero in probability
by tightness, its second in mean, and its third uniformly.
The expectation difference on the left is bounded by $2\norm G_\infty$,
so its absolute value, restricted to good exteriors, tends to zero in mean.

Let $\mathfrak F_K$ be the statistic $F_K$ evaluated on the rescaled
random block. Since $E_{K,M}$ is retained-exterior measurable, conditional
DLR gives
\[
\limsup_{K\to\infty}
\abs{\E_P\mathfrak F_K-\omega_{K,T_K}(F_K)}
\le 2\norm G_\infty P(g_*>M).
\]
The same estimate holds for $Q$, with exactly the same reference.
Subtract, let $M\to\infty$, and obtain
$\E_P\mathfrak F_K-\E_Q\mathfrak F_K\to0$.
The rescaling and Palm-identification lemmas of
Section~\ref{sec:prelim} then identify all finite consecutive-gap
distributions of $P^0$ and $Q^0$.
Point-shift invariance identifies every finite window of integer
gap indices. These gaps determine the entire rooted configuration.
Palm inversion at intensity one therefore gives $P=Q$.
\end{proof}

\subsection{Proof for inverse powers}
\begin{proof}[Proof of Theorem~\ref{thm:riesz}]
All statements of Section~\ref{sec:prelim} follow from
\eqref{hyp:linear}--\eqref{hyp:sublinear}. In
Section~\ref{sec:forces}, use the exponent $\alpha=s$.
Its weighted discrepancy condition is automatic, because
\begin{equation}\label{end:weighted-auto}
\int_1^\infty v_P(t)t^{-2s-2}\dd t
\le C_P\int_1^\infty t^{-2s-1}\dd t<\infty.
\end{equation}
Consequently the conditional cost $\mathcal D_{K,s}$ of the
difference of the scaled actual and lattice exterior forces is tight
on the good events. The same is true for $Q$.

Write $c=L/\ell$ and $\tau=c^s$. The exact homogeneity
\[
V_s(r/c)=c^sV_s(r)
\]
means that the rescaled conditional DLR density has Hamiltonian
$\beta\tau H_{V_s,(0,L)}^{c(\zeta-a),K}$, up to an additive constant.
Compare it first with $\omega_{K,T}^{(\tau)}$, defined by
\eqref{def:omega-tau}. The internal pair interactions and the charges
at $0,L$ have identical coefficients in the two densities.
Their logarithmic derivatives therefore cancel exactly. What remains is
the exterior-force difference and the quadratic confinement. Thus
\begin{equation}\label{end:riesz-score}
I(\mu_\zeta\mid\omega_{K,T}^{(\tau)})
\le C_{\beta,s}\tau^2
\left(\mathcal D_{K,s}+\frac{J_K}{T^2}\right).
\end{equation}
For every fixed retained exterior, the residual exterior particles are
a positive distance from the endpoints. Their relative forces are
bounded on $[0,L]$ and their tails converge uniformly there.
The score in \eqref{end:riesz-score} is bounded on $\Omega_K$.
In particular both the relative entropy and Fisher information are
finite; no inverse-gap moment estimate has been used.

On $E_{K,M}$, eventually $c\in[1/2,2]$ and hence
$\tau\in[2^{-s},2^s]$. Apply the static gap inequality with inverse
temperature $\beta\tau$ and confinement parameter $T/\tau$.
The constants are uniform on this compact interval, so
\[
\abs{\mu_\zeta(F_K)-\omega_{K,T}^{(\tau)}(F_K)}^2
\le C_G\left\{\frac{T}{K}\mathcal D_{K,s}
+\frac{J_K}{KT}\right\}.
\]
Corollary~\ref{sta:temperature-gap} gives
\[
\abs{\omega_{K,T}^{(\tau)}(F_K)-\omega_{K,T}(F_K)}^2
\le C_G(\tau-1)^2.
\]
Combining these inequalities proves \eqref{end:abstract} with
$\mathcal A_K=\mathcal D_{K,s}$ and $e_K=c^s-1$.
On good boundaries,
\[
|c^s-1|\le C_s|c-1|\le C_sr_KK^{-1/2}.
\]
The displacement estimate gives the remaining hypothesis of
Lemma~\ref{end:criterion}, which proves the theorem.
\end{proof}

\begin{remark}
It is important to scale the \emph{whole} reference exponent by
$\tau$. The confinement then has strength $\tau/T$, which stays
comparable to $1/T$. The temperature comparison applies directly to
this family and avoids all singular internal-force estimates.
\end{remark}

\subsection{The dilation error for smooth perturbations}
In this subsection $V=-a_0\log r+W$ satisfies
\eqref{hyp:convex}. The conditional law $\mu_\zeta$ after dilation has
interaction $r\mapsto V(r/c)$; the deterministic reference uses $V(r)$.
The logarithmic interior terms cancel because
$-a_0\log(r/c)=-a_0\log r+a_0\log c$.

\begin{lemma}[The added internal score]\label{end:internal-W}
Assume \eqref{hyp:third}. Define
\begin{align*}
A_c(r)&=c^{-1}W'(r/c)-W'(r),\\
R_i(u)&=\sum_{j\ne i}\sgn(u_i-u_j)A_c(|u_i-u_j|),\\
\mathcal R_K(\zeta)&=\E_{\mu_\zeta}\sum_{i=1}^K|R_i(u)|^2.
\end{align*}
For every fixed $M$,
\[
\E_P[\1_{E_{K,M}}\mathcal R_K]\longrightarrow0.
\]
\end{lemma}

\begin{proof}
For $c\in[1/2,2]$,
\[
\frac{\partial}{\partial c}\bigl[c^{-1}W'(r/c)\bigr]
=-c^{-2}W'(r/c)-rc^{-3}W''(r/c).
\]
The derivatives of $W$ are bounded and vanish beyond $R_0$.
Integration in $c$ therefore gives
\[
|A_c(r)|\le C_W|c-1|\1_{\{r\le2R_0\}}.
\]
On $E_{K,M}$ the physical interior points belong to $(0,3K]$ for
large $K$. A scaled separation at most $2R_0$ corresponds to a
physical separation at most $4R_0$. It follows that
\begin{equation}\label{end:internal-physical}
\sum_i|R_i(u)|^2
\le C_W|c-1|^2
\sum_{x\in\gamma\cap(0,3K]}
N_{[x-4R_0,x+4R_0]}(\gamma)^2.
\end{equation}

The moment assumption implies
$\E_{P^0}N_{[-4R_0,4R_0]}^2<\infty$. Indeed, Campbell's formula gives
\[
\begin{split}
\E_{P^0}N_{[-4R_0,4R_0]}^2
&=\E_P\sum_{x\in\gamma\cap(0,1]}
N_{[x-4R_0,x+4R_0]}^2\\
&\le \E_PN_{[-4R_0,1+4R_0]}^3<\infty.
\end{split}
\]
The last assertion follows by covering the fixed interval by finitely
many unit intervals and using stationarity.
Another application of Campbell yields
\[
\E_P\sum_{x\in\gamma\cap(0,3K]}
N_{[x-4R_0,x+4R_0]}^2
=3K\,\E_{P^0}N_{[-4R_0,4R_0]}^2\le CK.
\]
Since $|c-1|^2\le Cr_K^2/K$ on $E_{K,M}$, taking expectations
in \eqref{end:internal-physical} gives an upper bound $Cr_K^2$.
The good event is exterior measurable, so conditioning gives the same
bound for $\E_P[\1_{E_{K,M}}\mathcal R_K]$.
\end{proof}

\begin{lemma}[The added exterior score]\label{end:exterior-W}
Let $\mathcal W_K(\zeta)$ be the squared conditional force cost of the
$W$-part of the actual scaled exterior interaction minus the $W$-part
of the lattice exterior interaction, including their endpoint terms.
Then $\1_{E_{K,M}}\mathcal W_K$ is tight for every fixed $M$.
\end{lemma}

\begin{proof}
The scaled actual force has summands
$c^{-1}\sgn(u-p)W'(|u-p|/c)$, whereas the lattice force has summands
$\sgn(u-p)W'(|u-p|)$. On good boundaries their magnitudes are bounded
by a deterministic constant and they vanish unless the exterior charge
lies within scaled distance $2R_0$ of $u$.
Consequently the left contribution vanishes for $u>2R_0$, and the
right one vanishes for $L-u>2R_0$. Let
\[
Y_+=1+N_{[a-4R_0,a)}(\gamma),\qquad
Y_-=1+N_{(b,b+4R_0]}(\gamma).
\]
These variables are exterior measurable. The constants include the
common endpoint and the finitely many lattice charges that can contribute.
Writing $X_\pm$ for the conditional interior count envelopes,
\[
\mathcal W_K\le C_{W,R_0}
\bigl(Y_+^2X_++Y_-^2X_-\bigr).
\]
For example, the number of left interior particles affected is at most
$(1+2R_0)C_{\mathrm{int}}^+$.
On $\{g_*\le M\}$ the laws of both endpoint environments
are dominated by $M P^0$, uniformly in $K$, so $Y_\pm$ are tight.
The restricted first-moment bounds for $X_\pm$ give their tightness.
Products of nonnegative tight random variables are tight.
\end{proof}

\begin{proof}[Proof of Theorem~\ref{thm:perturb}]
The static inequality of Section~\ref{sec:static} applies to the
reference with the full potential $V$, because \eqref{hyp:convex}
ensures all pair Hessians are nonnegative and the necessary compact
gap ranges have positive curvature.

Let $\mathcal D_{K,0}$ be the logarithmic exterior-force cost from
Section~\ref{sec:forces}; it is tight on good boundaries by
\eqref{hyp:integrated}. Cancellation of the interior logarithms and
endpoint logarithms, followed by the two preceding lemmas, gives
\begin{equation}\label{end:W-score}
I(\mu_\zeta\mid\omega_{K,T})
\le C_{\beta,a_0,W}
\left(\mathcal D_{K,0}+\mathcal W_K+\mathcal R_K+\frac{J_K}{T^2}\right).
\end{equation}
For fixed exterior, the residual logarithmic exterior force is bounded
on the closed interval, and all $W$ derivatives appearing in the score
are bounded. Hence the entropy and Fisher information in
\eqref{end:W-score} are finite for every $\beta>0$.

The static Fisher estimate proves \eqref{end:abstract} with
$\mathcal A_K=\mathcal D_{K,0}+\mathcal W_K+\mathcal R_K$ and $e_K=0$.
The first two terms are tight, and the third tends to zero in
restricted mean. The displacement cost is $o(K^2)$ in restricted mean.
Lemma~\ref{end:criterion} proves uniqueness.
\end{proof}

\begin{remark}
The third count moment is a sufficient hypothesis for the particular
internal-score estimate used here; we do not assert it is necessary.
When $W=0$, the two added costs vanish and this moment assumption can
be removed, recovering the discrepancy comparison theorem for the
unperturbed logarithm.
\end{remark}


\section{Scope of the extensions and higher-dimensional obstructions}
\label{sec:limits}

\subsection{Collision regularity for more general perturbations}
Positive curvature alone does not ensure that a direct comparison
after dilation has finite Fisher information. Suppose, near $r=0$,
\[
V(r)=-a_0\log r+b r^\eta,\qquad
0<\eta<\tfrac12,\quad b\ne0.
\]
For $c\ne1$ the added interior-force defect is
$b\eta(c^{-\eta}-1)r^{\eta-1}$. On a region where just one pair
collides and the remaining separations stay bounded away from zero,
the Gibbs density is comparable to $r^{\beta a_0}$. The squared
defect therefore has the collision integral
\[
\int_0^\delta r^{\beta a_0+2\eta-2}\dd r,
\]
which diverges if $\beta a_0+2\eta\le1$.
One can multiply $r^\eta$ by a smooth cutoff and choose $|b|$
small enough that the full $V''$ remains positive: near zero the
$a_0r^{-2}$ term dominates, and away from zero the cutoff contribution
is bounded on a compact annulus. Thus this obstruction can occur
within convex interactions. It concerns the direct Fisher proof,
not the truth of uniqueness.

\subsection{Why a finite-energy theorem needs another input}
Theorems~\ref{thm:riesz} and~\ref{thm:perturb} assume the count bounds.
They do not deduce them from finite energy. For $0<s<1$, the
one-dimensional Riesz Fourier kernel is a positive multiple of
$|k|^{s-1}$. For example, with the Fourier convention
$\widehat f(k)=\int e^{-ikx}f(x)\dd x$, the Gaussian integral
representation gives, distributionally,
\[
|x|^{-s}=\frac1{\Gamma(s/2)}
\int_0^\infty t^{s/2-1}e^{-tx^2}\dd t,\qquad
\widehat{|\,\cdot\,|^{-s}}(k)
=\frac{2^{1-s}\sqrt\pi\,\Gamma((1-s)/2)}{\Gamma(s/2)}
|k|^{s-1}.
\]
Consequently a flat covariance spectral density has finite
low-frequency Riesz cost:
\[
\int_{|k|<1}|k|^{s-1}\dd k=\frac2s<\infty.
\]
Smooth charge regularization controls the high-frequency part.
In contrast, the logarithmic weight $|k|^{-1}$ is not integrable
against a flat density near zero. The logarithmic energy argument
therefore cannot by itself prove sublinear variance for a Riesz state.
This observation does not assert that a Riesz DLR state has flat
spectrum; the DLR equations may provide additional information.
For $s\ge1$, the hypotheses of Theorem~\ref{thm:riesz} also remain to be
verified in whatever class of Gibbs states is under consideration.

\subsection{Negative tangential curvature}\label{lim:tangential}
\begin{proposition}\label{lim:hessian}
Let $d\ge2$ and $V\in C^2((0,\infty))$. For $z\ne0$, put
$r=|z|$ and $e=z/r$. Then
\begin{equation}\label{lim:radial-Hessian}
D^2[V(|z|)]=V''(r)e\otimes e+
\frac{V'(r)}r(I-e\otimes e).
\end{equation}
In particular the radial logarithmic and positive inverse-power
interactions do not satisfy the pointwise convexity needed in
Section~\ref{sec:static}. A fixed quadratic confinement does not remove
this failure near collisions.
\end{proposition}

\begin{proof}
Differentiate $\nabla[V(|z|)]=V'(r)e$ and use
$D e=(I-e\otimes e)/r$. The radial eigenvalue is $V''(r)$, while
each tangential eigenvalue is $V'(r)/r$. Thus the eigenvalues are
\[
\begin{array}{c|cc}
V(r)&\text{radial}&\text{tangential}\\ \hline
-\log r&r^{-2}&-r^{-2}\\[2pt]
r^{-s}&s(s+1)r^{-s-2}&-sr^{-s-2}
\end{array}
\]
for $s>0$.
Place two interior particles at separation $re$ and choose their
velocities to be $h,-h$, with $h\perp e$.
Their pair contribution to the Hamiltonian Hessian is
$4V'(r)|h|^2/r$, which tends to $-\infty$ as $r\downarrow0$.
Keep all other particles a positive distance from this pair and
assume the fixed exterior potential is $C^2$ near it.
All other Hessian contributions, including the quadratic confinement,
then remain bounded. This proves the assertion.
\end{proof}

The proposition includes Coulomb repulsion: $-\log r$ in dimension
two and $r^{2-d}$ in dimension $d\ge3$. It proves a limitation of
this convexity argument, rather than nonuniqueness of Gibbs states.
In a regime with distinct stationary phases, a global uniqueness
theorem would of course be false. Different translates of a phase
can become the same law after translation averaging, whereas distinct
orientational phases can retain different stationary averages.
Neither existence nor coexistence of such phases is proved here.

\subsection{Geometry, angular discrepancy, and scales}
The ordered one-dimensional construction also requires replacement:
there is no direct counterpart of its successor-gap coding and
two matched endpoint charges in a general free-particle gas.
Scalar centered-ball counts alone do not determine the exterior force.
For example, the signed charge
$\delta_{Re_1}-\delta_{-Re_1}$ has zero mass in every ball centered
at the origin but produces a nonzero force there when $V'(R)\ne0$.
An extension needs angular or field information.

The growth of boundary area alone is not an obstruction to the
abstract comparison. Suppose a different analytic argument produced
\[
|\mu(F_R)-\omega(F_R)|^2
\le C\left(\frac{T\mathcal D_R}{N_R}+
\frac{J_R}{N_RT}\right),\qquad N_R\asymp R^d.
\]
If $\mathcal D_R=O_{\Prob}(R^{d-1})$ and
$\E J_R=N_R\sigma_R^2$, the two scales would be $T/R$ and
$\sigma_R^2/T$. They can both vanish if
\[
\sigma_R^2=o(R),\qquad \sigma_R^2\ll T\ll R.
\]
For instance, $T=\sqrt{R\max\{1,\sigma_R^2\}}$ works.
This is a conditional scale calculation: the functional inequalities
and the surface-order force bound have not been proved here.

\subsection{A finite-field-energy obstruction above dimension two}
\begin{proposition}\label{lim:poisson}
In every dimension $d\ge3$ there is a stationary point process with
linear volume-order count variance admitting a stationary compatible
electric field of finite mean square after fixed-scale smooth charge
regularization.
\end{proposition}

\begin{proof}
Let $\Pi$ be the unit Poisson process and let $-\Delta G=\delta_0$.
Choose a smooth compactly supported probability density $\rho_\eta$,
and set $K_\eta=(\nabla G)*\rho_\eta$.
This vector field is locally bounded and satisfies
$|K_\eta(x)|=O(|x|^{1-d})$ at infinity, so
\[
\int_{\R^d}|K_\eta|^2<\infty
\]
because its tail is bounded by a constant times
$\int_1^\infty r^{1-d}\dd r$.
The compensated Poisson integral
\[
E(x)=\int K_\eta(x-y)(\Pi(\dd y)-\dd y)
\]
exists in $L^2$, is stationary, and has
\[
\E|E(0)|^2=\int_{\R^d}|K_\eta(y)|^2\dd y<\infty.
\]
Since $-\diver K_\eta=\rho_\eta$, testing against smooth compactly
supported functions gives
$-\diver E=\Pi*\rho_\eta-1$ in distributions.
However $\Var\Pi(B_R)=|B_R|$. Thus the stated field-energy property
does not force sublinear volume-order count variance.
\end{proof}

This is a counterexample to an energy-only implication. The Poisson
process is not asserted to satisfy Coulomb DLR, and fixed-scale
regularized energy is not identified here with a full renormalized
energy convention.

\begin{thebibliography}{DHLM}
\bibitem[Draft]{Draft}
\emph{Uniqueness of the stationary logarithmic DLR state}.
User-supplied working manuscript, file
\texttt{logarithmic\_dlr\_uniqueness.tex}, dated 9 September 2026.
The source of the comparison framework developed in this note.

\bibitem[BEY]{BEY}
P.~Bourgade, L.~Erd\H{o}s and H.-T.~Yau,
\emph{Universality of general $\beta$-ensembles}.
Duke Math.\ J.\ \textbf{163} (2014), no.~6, 1127--1190.
\href{https://arxiv.org/abs/1104.2272}{arXiv:1104.2272}.

\bibitem[DHLM]{DHLM}
D.~Dereudre, A.~Hardy, T.~Lebl\'e and M.~Ma\"ida,
\emph{DLR equations and rigidity for the Sine-beta process}.
Preprint version 2, 2019.
\href{https://arxiv.org/abs/1809.03989v2}{arXiv:1809.03989v2}.

\bibitem[FMW]{FMW}
M.~Fradelizi, M.~Madiman and L.~Wang,
\emph{Optimal concentration of information content for log-concave
densities}.
In \emph{High Dimensional Probability VII}, Progress in Probability
\textbf{71}, Springer, 2016.
Preprint version 2:
\href{https://arxiv.org/abs/1508.04093v2}{arXiv:1508.04093v2}.
\end{thebibliography}
\end{document}
